\documentclass[11pt,a4paper]{article}
\usepackage[margin=2.0cm]{geometry}
\usepackage[T1]{fontenc}
\usepackage{amsmath}
\usepackage{booktabs}
\usepackage{xcolor}
\usepackage{listings}
\usepackage[hidelinks]{hyperref}

\emergencystretch=2em
\definecolor{codebg}{gray}{0.95}
\lstset{language=C++, basicstyle=\ttfamily\footnotesize, keywordstyle=\color{blue!70!black},
  commentstyle=\color{green!50!black}\itshape, backgroundcolor=\color{codebg},
  frame=single, rulecolor=\color{black!20}, breaklines=true, showstringspaces=false,
  tabsize=2, columns=flexible}

\title{ATmega32A Module \\ \Large GPIO (Digital Input / Output)}
\author{Ali Sahafi}
\date{}

\begin{document}
\maketitle

\section{Theory}

GPIO (\emph{General Purpose Input/Output}) is the most basic peripheral: each
pin of the microcontroller can be a digital \textbf{output} (drive 0\,V or
5\,V) or a digital \textbf{input} (read whether the outside world applies
0\,V or 5\,V). The ATmega32A has four 8-bit ports: \textbf{PORTA},
\textbf{PORTB}, \textbf{PORTC} and \textbf{PORTD} --- 32 pins in total.

Each port is controlled by three registers ($x$ = A, B, C or D):

\begin{center}
\begin{tabular}{lll}
\toprule
\textbf{Register} & \textbf{Role} & \textbf{Bit meaning} \\
\midrule
\texttt{DDR$x$}  & Data Direction & 1 = output, 0 = input \\
\texttt{PORT$x$} & Output latch / pull-up & output: pin level; input: 1 = pull-up on \\
\texttt{PIN$x$}  & Input read & the actual voltage currently on the pin \\
\bottomrule
\end{tabular}
\end{center}

\paragraph{Internal pull-up.} An unconnected input pin \emph{floats}: it reads
random values. For buttons, enable the internal pull-up
(\texttt{INPUT\_PULLUP}): the pin then rests at \texttt{HIGH}, and the button
connects it to GND, so \textbf{pressed = LOW}. This ``active-low'' logic is
the standard way to wire buttons.

\section{API reference}

The driver is beginner-friendly and auto-corrects register names: you can pass
\texttt{PORTx}, \texttt{DDRx}, or \texttt{PINx} interchangeably to any function
(\texttt{setDirection}, \texttt{write}, \texttt{read}, \texttt{toggle}). The driver
automatically maps to the correct hardware register under the hood, skipping register
confusion.

\begin{center}
\small
\begin{tabular}{ll}
\toprule
\textbf{Method} & \textbf{Description} \\
\midrule
\texttt{GPIO.setDirection(PORTx, pin, OUTPUT)} & Set single pin (\texttt{INPUT}, \texttt{OUTPUT}, \texttt{INPUT\_PULLUP}) \\
\texttt{GPIO.setDirection(PORTx, ALL, INPUT\_PULLUP)} & Set whole port direction \\
\texttt{GPIO.write(PORTx, pin, HIGH/LOW)} & Write single pin \\
\texttt{GPIO.write(PORTx, ALL, 0xF0)} & Write raw byte to whole port \\
\texttt{GPIO.write(PORTx, 0xF0)} & Shorthand whole-port write \\
\texttt{GPIO.read(PORTx, pin)} & Read single pin --- returns \texttt{HIGH} or \texttt{LOW} \\
\texttt{GPIO.read(PORTx, ALL)} & Read whole port --- returns 0--255 \\
\texttt{GPIO.read(PORTx)} & Shorthand whole-port read \\
\texttt{GPIO.toggle(PORTx, pin)} & Toggle single pin \\
\texttt{GPIO.toggle(PORTx, ALL)} & Toggle whole port \\
\bottomrule
\end{tabular}
\end{center}

\textbf{Constants:} \texttt{INPUT}, \texttt{OUTPUT}, \texttt{INPUT\_PULLUP},
\texttt{HIGH}, \texttt{LOW}, \texttt{ALL} (see Module 0).

\newpage
\section{Example: Button-controlled LED}

\paragraph{Board Hardware.} All components are already integrated on the development board:
\begin{itemize}
  \item \textbf{8 LEDs (D0--D7) on PORTB (PB0--PB7):} Wired active-low (cathodes to MCU pins). Driving a pin \texttt{LOW} turns the LED ON; driving it \texttt{HIGH} turns it OFF.
  \item \textbf{8 DIP Switches (S0--S7) on PORTC (PC0--PC7):} Connected directly to GND when closed. There are \textbf{no external pull-up resistors}, so \texttt{INPUT\_PULLUP} must be enabled in software.
  \item \textbf{2 Push Buttons on PD2 and PD6:} Button S11 is connected to \textbf{PD2} and button S10 is connected to \textbf{PD6}. Pressing either button connects the pin to GND (\texttt{LOW}).
\end{itemize}

\paragraph{Build \& flash.} \texttt{make gpio}

\lstinputlisting{example.cpp}

\newpage
\section{Laboratory Tasks}

\subsection*{Task 1: Reading Switches and Driving LEDs (20 min)}
\begin{itemize}
  \item \textbf{Task 1.1:} Configure \texttt{PORTC} as \texttt{INPUT\_PULLUP} and \texttt{PORTB} as \texttt{OUTPUT}. Write a loop that continuously reads the 8 DIP switches on \texttt{PORTC} and displays their status directly on the 8 LEDs on \texttt{PORTB}.
  \item \textbf{Task 1.2:} Modify the configuration of \texttt{PORTC} to plain \texttt{INPUT} (disabling the internal pull-up resistors). Observe and explain what happens when a switch is left open (floating input).
\end{itemize}

\subsection*{Task 2: 4-Bit Binary to Seven-Segment Decoder (20 min)}
Develop a function to capture a 4-bit binary input from the switches and display the corresponding decimal digit on the seven-segment display:
\begin{itemize}
  \item \textbf{Input Capture:} Read the lower 4 bits of the DIP switches on \texttt{PORTC} (\texttt{PC0}--\texttt{PC3}) using bitwise masking (\texttt{value \& 0x0F}).
  \item \textbf{Value Mapping:} Convert the 4-bit binary value (0--9) into the corresponding segment pattern for the active-low common-anode display (refer to the segment encoding table from the lecture slides).
  \item \textbf{Output to Display:} Send the mapped segment byte to \texttt{PORTB} to illuminate the digit on the display.
  \item \textbf{Example:} If the switches are set to binary \texttt{0101}, the seven-segment display should show the digit ``5''.
\end{itemize}

\subsection*{Task 3: Bidirectional Running Light with Push-Button Control (25 min)}
Apply shift operators (\texttt{<<}, \texttt{>>}), bitwise logic (\texttt{\&}, \texttt{\textasciitilde}), and \texttt{\_delay\_ms()} to create an animated running light on \texttt{PORTB}:
\begin{itemize}
  \item \textbf{Task 3.1 (Shift Animation):} Use a single active bit (e.g., \texttt{0x01}) and shift it left and right across the 8 LEDs on \texttt{PORTB} (\texttt{PB0} $\leftrightarrow$ \texttt{PB7}) using shift operators (\texttt{<<} and \texttt{>>}) and a delay of \texttt{\_delay\_ms(100)}. Invert the pattern (\texttt{\textasciitilde pattern}) when writing to \texttt{PORTB} since the on-board LEDs are active-low.
  \item \textbf{Task 3.2 (Push-Button Control):} Integrate the two on-board push buttons on \texttt{PORTD} to control the animation:
    \begin{itemize}
      \item Button S11 on \textbf{PD2}: Direct the running light to shift to the left.
      \item Button S10 on \textbf{PD6}: Direct the running light to shift to the right.
    \end{itemize}
  \item \textbf{Task 3.3 (Variable Speed via Switches):} Read the lower 2 bits of the \texttt{PORTC} DIP switches using a bitwise mask (\texttt{\& 0x03}) and use an arithmetic expression to dynamically scale the delay duration (e.g., $50\,\text{ms} \times (\text{setting} + 1)$ for 50\,ms, 100\,ms, 150\,ms, and 200\,ms).
\end{itemize}

\end{document}
